% Options for packages loaded elsewhere
% Options for packages loaded elsewhere
\PassOptionsToPackage{unicode}{hyperref}
\PassOptionsToPackage{hyphens}{url}
\PassOptionsToPackage{dvipsnames,svgnames,x11names}{xcolor}
%
\documentclass[
  turkish,
  12pt,
  a4paper,
]{report}
\usepackage{xcolor}
\usepackage[top=2cm,bottom=2cm,left=2.5cm,right=2.5cm]{geometry}
\usepackage{amsmath,amssymb}
\setcounter{secnumdepth}{5}
\usepackage{iftex}
\ifPDFTeX
  \usepackage[T1]{fontenc}
  \usepackage[utf8]{inputenc}
  \usepackage{textcomp} % provide euro and other symbols
\else % if luatex or xetex
  \usepackage{unicode-math} % this also loads fontspec
  \defaultfontfeatures{Scale=MatchLowercase}
  \defaultfontfeatures[\rmfamily]{Ligatures=TeX,Scale=1}
\fi
\usepackage{lmodern}
\ifPDFTeX\else
  % xetex/luatex font selection
  \setmainfont[]{TeX Gyre Termes}
\fi
% Use upquote if available, for straight quotes in verbatim environments
\IfFileExists{upquote.sty}{\usepackage{upquote}}{}
\IfFileExists{microtype.sty}{% use microtype if available
  \usepackage[]{microtype}
  \UseMicrotypeSet[protrusion]{basicmath} % disable protrusion for tt fonts
}{}
\makeatletter
\@ifundefined{KOMAClassName}{% if non-KOMA class
  \IfFileExists{parskip.sty}{%
    \usepackage{parskip}
  }{% else
    \setlength{\parindent}{0pt}
    \setlength{\parskip}{6pt plus 2pt minus 1pt}}
}{% if KOMA class
  \KOMAoptions{parskip=half}}
\makeatother
% Make \paragraph and \subparagraph free-standing
\makeatletter
\ifx\paragraph\undefined\else
  \let\oldparagraph\paragraph
  \renewcommand{\paragraph}{
    \@ifstar
      \xxxParagraphStar
      \xxxParagraphNoStar
  }
  \newcommand{\xxxParagraphStar}[1]{\oldparagraph*{#1}\mbox{}}
  \newcommand{\xxxParagraphNoStar}[1]{\oldparagraph{#1}\mbox{}}
\fi
\ifx\subparagraph\undefined\else
  \let\oldsubparagraph\subparagraph
  \renewcommand{\subparagraph}{
    \@ifstar
      \xxxSubParagraphStar
      \xxxSubParagraphNoStar
  }
  \newcommand{\xxxSubParagraphStar}[1]{\oldsubparagraph*{#1}\mbox{}}
  \newcommand{\xxxSubParagraphNoStar}[1]{\oldsubparagraph{#1}\mbox{}}
\fi
\makeatother


\usepackage{longtable,booktabs,array}
\newcounter{none} % for unnumbered tables
\usepackage{calc} % for calculating minipage widths
% Correct order of tables after \paragraph or \subparagraph
\usepackage{etoolbox}
\makeatletter
\patchcmd\longtable{\par}{\if@noskipsec\mbox{}\fi\par}{}{}
\makeatother
% Allow footnotes in longtable head/foot
\IfFileExists{footnotehyper.sty}{\usepackage{footnotehyper}}{\usepackage{footnote}}
\makesavenoteenv{longtable}
\usepackage{graphicx}
\makeatletter
\newsavebox\pandoc@box
\newcommand*\pandocbounded[1]{% scales image to fit in text height/width
  \sbox\pandoc@box{#1}%
  \Gscale@div\@tempa{\textheight}{\dimexpr\ht\pandoc@box+\dp\pandoc@box\relax}%
  \Gscale@div\@tempb{\linewidth}{\wd\pandoc@box}%
  \ifdim\@tempb\p@<\@tempa\p@\let\@tempa\@tempb\fi% select the smaller of both
  \ifdim\@tempa\p@<\p@\scalebox{\@tempa}{\usebox\pandoc@box}%
  \else\usebox{\pandoc@box}%
  \fi%
}
% Set default figure placement to htbp
\def\fps@figure{htbp}
\makeatother



\ifLuaTeX
\usepackage[bidi=basic,shorthands=off]{babel}
\else
\usepackage[bidi=default,shorthands=off]{babel}
\fi
\ifPDFTeX
\else
\babelfont{rm}[]{TeX Gyre Termes}
\fi
\ifLuaTeX
  \usepackage{selnolig} % disable illegal ligatures
\fi


\setlength{\emergencystretch}{3em} % prevent overfull lines

\providecommand{\tightlist}{%
  \setlength{\itemsep}{0pt}\setlength{\parskip}{0pt}}



 


% --- Marmara §1.2: 1,5 satir araligi + paragraf 6 nk + girinti yok ---
\usepackage{setspace}
\onehalfspacing
\usepackage{parskip}
\setlength{\parindent}{0pt}
\setlength{\parskip}{6pt plus 1pt}
% --- §1.2: iki yana yasli, satir sonunda kelime bolme yok ---
\usepackage[none]{hyphenat}
\sloppy
\hyphenpenalty=10000
\exhyphenpenalty=10000
% --- Turkce heceleme/karakter (polyglossia XeLaTeX icin) ---
\usepackage{polyglossia}
\setmainlanguage{turkish}
% --- §5 dizin basligi ezimi include-before-body'de yapilir; Quarto
% kendi \AtBeginDocument ezimini enjekte ettiginden (ve o bizimkinden
% sonra calistigindan) yeniden-adlandirma \listoffigures'tan hemen
% once, govde basinda uygulanir (bkz. include-before-body). ---
% --- §1.9: sayfa numarasi alt-orta ---
\usepackage{fancyhdr}
\pagestyle{fancy}
\fancyhf{}
\fancyfoot[C]{\thepage}
\renewcommand{\headrulewidth}{0pt}
\renewcommand{\footrulewidth}{0pt}
\fancypagestyle{plain}{%
  \fancyhf{}%
  \fancyfoot[C]{\thepage}%
  \renewcommand{\headrulewidth}{0pt}%
}
% --- §1.3: bolum basligi bold 14 pt, alt basliklar bold 12 pt ---
\usepackage{titlesec}
\titleformat{\chapter}{\normalfont\fontsize{14}{17}\bfseries}{\thechapter.}{1em}{}
\titleformat{\section}{\normalfont\fontsize{12}{15}\bfseries}{\thesection}{1em}{}
\titleformat{\subsection}{\normalfont\fontsize{12}{15}\bfseries}{\thesubsection}{1em}{}
\titleformat{\subsubsection}{\normalfont\fontsize{12}{15}\bfseries}{\thesubsubsection}{1em}{}
% bolum basligi sayfanin tepesinden baslasin (buyuk ust bosluk olmasin)
\titlespacing*{\chapter}{0pt}{0pt}{1\baselineskip}
% --- §1.9: kapak/onay/icindekiler/dizin sayfalari numarasiz baslar;
% Kisaltmalar'dan itibaren Romen, Ozet'ten itibaren Arap rakami
% (gecisler chapters/00b ve 00c icindeki ham LaTeX bloklariyla). ---
\pagenumbering{gobble}
\makeatletter
\@ifpackageloaded{caption}{}{\usepackage{caption}}
\AtBeginDocument{%
\ifdefined\contentsname
  \renewcommand*\contentsname{İçindekiler}
\else
  \newcommand\contentsname{İçindekiler}
\fi
\ifdefined\listfigurename
  \renewcommand*\listfigurename{Şekil Listesi}
\else
  \newcommand\listfigurename{Şekil Listesi}
\fi
\ifdefined\listtablename
  \renewcommand*\listtablename{Tablo Listesi}
\else
  \newcommand\listtablename{Tablo Listesi}
\fi
\ifdefined\figurename
  \renewcommand*\figurename{Şekil}
\else
  \newcommand\figurename{Şekil}
\fi
\ifdefined\tablename
  \renewcommand*\tablename{Tablo}
\else
  \newcommand\tablename{Tablo}
\fi
}
\@ifpackageloaded{float}{}{\usepackage{float}}
\floatstyle{ruled}
\@ifundefined{c@chapter}{\newfloat{codelisting}{h}{lop}}{\newfloat{codelisting}{h}{lop}[chapter]}
\floatname{codelisting}{Listeleme}
\newcommand*\listoflistings{\listof{codelisting}{İlan Listesi}}
\captionsetup{labelsep=period}
\makeatother
\makeatletter
\makeatother
\makeatletter
\@ifpackageloaded{caption}{}{\usepackage{caption}}
\@ifpackageloaded{subcaption}{}{\usepackage{subcaption}}
\makeatother
\usepackage{bookmark}
\IfFileExists{xurl.sty}{\usepackage{xurl}}{} % add URL line breaks if available
\urlstyle{same}
\hypersetup{
  pdflang={tr},
  colorlinks=true,
  linkcolor={black},
  filecolor={Maroon},
  citecolor={black},
  urlcolor={black},
  pdfcreator={LaTeX via pandoc}}


\author{}
\date{}
\begin{document}

% §5: Quarto'nun Turkce LOF/LOT ad ezimi \AtBeginDocument ile bizim
% header ezimimizi geciyor; govde basinda (tum \AtBeginDocument
% hook'lari calistiktan SONRA) yeniden adlandirarak sabitliyoruz.
\renewcommand{\listfigurename}{ŞEKİLLER DİZİNİ}
\renewcommand{\listtablename}{TABLOLAR DİZİNİ}

\renewcommand*\contentsname{İÇİNDEKİLER}
{
\hypersetup{linkcolor=black}
\setcounter{tocdepth}{2}
\tableofcontents
}
\listoffigures
\listoftables

\chapter{Spec --- Kapsamlı Tez Kontrol Checklisti + Doğrulama
Script'i}\label{spec-kapsamlux131-tez-kontrol-checklisti-doux11frulama-scripti}

\section{Problem Tanımı}\label{problem-tanux131mux131}

Repo, olgun bir kalite-kontrol katmanına sahip:
\texttt{tez-yazim/04\_kalite-kontrol/} altında 6-kapı (Kapı 0--5)
sertifikasyon playbook'u, sertifika şablonu ve eksen bazlı
alt-checklistler (\texttt{format-kontrol-listesi.md},
\texttt{kanit-ve-gizlilik-kontrol-listesi.md},
\texttt{ai-mcp-kullanim-kontrol-listesi.md},
\texttt{turkce-bilimsel-yazim-denetimi.md},
\texttt{render-bagimliliklari.md}) mevcut. Ayrıca \texttt{scripts/util/}
altında çalışan doğrulama araçları var (\texttt{bib\_hygiene.py},
\texttt{karma\_ledger\_check.py}, \texttt{claim\_certification.py},
\texttt{csr\_numeric\_trace\_audit.py},
\texttt{csr\_causal\_label\_audit.py}, \texttt{csr\_block\_review.py},
\texttt{tr\_corpus\_audit.py}) ve slash komutları (\texttt{/tez-oturum},
\texttt{/bolum-sertifika}, \texttt{/referans-kapisi},
\texttt{/tez-literatur}, \texttt{/tez-dogrulama}).

Ancak \textbf{tek, uçtan uca, tüm tezi kapsayan bir ``master kontrol
checklisti'' yok}. Mevcut checklistler eksen-parçalıdır ve her madde
için (a) somut kontrol mekanizması, (b) çalıştırılabilir komut/eşik ve
(c) oto/manuel etiketi tek yerde toplanmış değildir. Bunun sonucu:

\begin{enumerate}
\def\labelenumi{\arabic{enumi}.}
\tightlist
\item
  \textbf{Dağınık doğrulama:} Bir bölümü veya tüm tezi teslim-öncesi
  kontrol etmek için hangi aracın hangi sırada koşacağı ve neyin PASS
  sayılacağı tek bir izlenebilir belgede yok.
\item
  \textbf{Otomasyon boşluğu:} Araçlar tek tek elle koşuluyor; hepsini
  sırayla koşup birleşik PASS/FAIL raporu üreten tek bir orkestratör
  script yok.
\item
  \textbf{Tez-düzeyi eksenler eksik:} Format bütünlüğü (Marmara §1),
  render/crossref bütünlüğü, targets pipeline tazeliği, PII/gizlilik
  sınırları, sayısal-iddia izlenebilirliği gibi eksenler
  bölüm-sertifikasında kısmen var ama tez-düzeyi toplu kapıda
  birleştirilmemiş.
\end{enumerate}

\section{Kapsam}\label{kapsam}

Kullanıcı kararları: - \textbf{Kapsam:} İki katmanlı --- hem
\textbf{bölüm-düzeyi} (her \texttt{chapters/*.qmd} için tekrar
koşulabilir) hem \textbf{tez-düzeyi} (kapanış/teslim öncesi bütün tez)
checklist, tek belgede. - \textbf{Mevcut sistemle ilişki:} Mevcut 6-kapı
checklistlerini \textbf{genişlet} --- eksik eksenleri var olan yapıya
ekle, master checklist onları birleştirip bağlar (sıfırdan paralel
sistem değil). - \textbf{Madde mekanizması:} Her madde için \textbf{hem}
çalıştırılabilir komut + beklenen çıktı + PASS eşiği \textbf{hem}
Otomatik/Manuel etiketi + kanıt konumu. - \textbf{Çıktı:}
\textbf{Markdown master checklist + Python doğrulama script'i}
(\texttt{scripts/util/} deseninde). - \textbf{Dil:} Türkçe, çok
kapsamlı.

\subsection{Script davranışı
(kritik)}\label{script-davranux131ux15fux131-kritik}

Doğrulama script'i \textbf{tüm otomatikleştirilebilir araçları sırayla
koşar ve birleşik rapor üretir}; hiçbir dosyayı değiştirmez (salt-okuma
denetim). Her madde için PASS/FAIL/SKIP + kısa kanıt satırı basar.
Manuel maddeler ``MANUEL --- elle kontrol gerekli'' olarak işaretlenir
(script bunları FAIL saymaz, raporda ayrı listeler). Script bir özet
exit kodu döndürür (0 = tüm otomatik maddeler PASS; \textgreater0 = en
az bir otomatik HARD fail) ki hem elde hem CI/pre-teslim akışında
kullanılabilsin.

\subsection{Kapsam dışı}\label{kapsam-dux131ux15fux131}

\begin{itemize}
\tightlist
\item
  Mevcut chapter içeriğini düzeltmek/yeniden yazmak (checklist tespit
  eder, düzeltme ayrı iştir).
\item
  Yeni bilimsel analiz veya pipeline hedefi eklemek.
\item
  CI altyapısı (GitHub Actions vb.) kurmak --- script CI-uyumlu yazılır
  ama workflow tanımı bu iş kapsamında değil.
\end{itemize}

\section{Gereksinimler}\label{gereksinimler}

\subsection{R1 --- Master checklist
belgesi}\label{r1-master-checklist-belgesi}

\begin{itemize}
\tightlist
\item
  Yeni dosya:
  \texttt{tez-yazim/04\_kalite-kontrol/tez-kontrol-checklisti.md}.
\item
  Türkçe, çok kapsamlı; iki ana bölüm: \textbf{(A) Bölüm-düzeyi kapı
  seti} ve \textbf{(B) Tez-düzeyi kapanış seti}.
\item
  Belge, mevcut 6-kapı playbook'unu ve alt-checklistleri \textbf{kanonik
  otorite} olarak referanslar; onların yerine geçmez, üstlerinde
  birleştirici indeks olur.
\item
  Her checklist maddesi standart bir satır/blok şemasıyla yazılır (bkz.
  R3).
\end{itemize}

\subsection{R2 --- Eksen kapsamı (checklist neyi kontrol
eder)}\label{r2-eksen-kapsamux131-checklist-neyi-kontrol-eder}

Master checklist en az şu eksenleri madde madde kapsar: 1.
\textbf{Kapsam \& Gizlilik (Kapı 0):} PII sınırları (\texttt{ad/soyad}
düşürme,
\texttt{data/raw\textbar{}\ \ \ \ identified\textbar{}cleaned\textbar{}backup}
git-dışı), credential/\texttt{.lock} erişimi,
\texttt{.claude/settings.json} deny uyumu. 2. \textbf{Kanıt \& Literatür
(Kapı 1--2):} kaynaksız-sayı taraması, atıf künye bütünlüğü, ledger
mutabakatı, full-text/Zotero izi. 3. \textbf{Bölüm metni \& kılavuz
uyumu (Kapı 3):} Marmara format kontratı §1 (kenar boşluk, satır
aralığı, font, başlık düzeni, tablo/şekil başlık biçimi), bölüm sırası,
başlık numaralandırma (çift-numara yok). 4. \textbf{Türkçe imla, akış,
mantık (Kapı 4):} Türkçe bilimsel yazım denetimi, ondalık virgül,
retorik/insan-Türkçesi kontrolleri. 5. \textbf{AI-reliability \& teknik
doğrulama (Kapı 5):} iki-kol hook derlemesi, render (quarto),
crossref/atıf bütünlüğü, sayısal-iddia → CSV izi, nedensel-etiket
denetimi. 6. \textbf{Pipeline \& üretilebilirlik (tez-düzeyi):}
\texttt{renv::status()} temiz, \texttt{targets::tar\_make()} tazeliği
(outdated hedef yok), üretilmiş tablo/şekil artefaktlarının varlığı. 7.
\textbf{Render \& çıktı bütünlüğü (tez-düzeyi):}
\texttt{quarto\ render\ thesis.qmd} exit 0; HTML/PDF'te 0 çözülmemiş
\texttt{?@}; 0 kırık atıf; şekil gömme; PDF Marmara format ölçüleri (A4,
kenar boşluk, font). 8. \textbf{Ön/arka bölümler \& teslim
(tez-düzeyi):} özet/summary kelime sınırı + anahtar sözcükler,
kısaltmalar, içindekiler/şekil/tablo listeleri, 00a resmi alanlar, 06
özgeçmiş, 07 ek yer-tutucuları, referans listesi.

\subsection{R3 --- Madde şeması (her checklist maddesi
için)}\label{r3-madde-ux15femasux131-her-checklist-maddesi-iuxe7in}

Her madde şu alanları içerir: - \textbf{ID:} kısa benzersiz kod (ör.
\texttt{K0-PII-01}). - \textbf{Kontrol maddesi:} ne doğrulanıyor
(Türkçe, tek cümle). - \textbf{Mekanizma:} nasıl kontrol edilir (araç
adı / gözle bakış yöntemi). - \textbf{Komut:} çalıştırılabilir komut
satırı (otomatik maddelerde) veya ``MANUEL'' (gözle kontrol
maddelerinde). - \textbf{PASS ölçütü:} beklenen çıktı / eşik (ör. ``exit
0'', ``HARD=0'', ``0 çözülmemiş crossref''). - \textbf{Tür:}
\texttt{OTOMATIK} \textbar{} \texttt{MANUEL} \textbar{} \texttt{YARI}
(araç + gözle teyit). - \textbf{Kanıt konumu:} çıktının/kanıtın nereye
yazıldığı (rapor, sertifika, log). - \textbf{Kapı/Kaynak:} hangi kapıya
(0--5) ve hangi kanonik kural dosyasına bağlı.

\subsection{R4 --- Doğrulama script'i}\label{r4-doux11frulama-scripti}

\begin{itemize}
\tightlist
\item
  Yeni dosya: \texttt{scripts/util/tez\_checklist\_verify.py} (mevcut
  util Python desenine uygun: \texttt{PYTHONDONTWRITEBYTECODE},
  salt-okuma, satır-düzeyi veri sızdırmaz).
\item
  Script, R2'deki \textbf{otomatikleştirilebilir} maddeleri sırayla
  koşar: mevcut araçları (\texttt{bib\_hygiene.py},
  \texttt{karma\_ledger\_check.py}, \texttt{claim\_certification.py},
  \texttt{csr\_numeric\_trace\_audit.py},
  \texttt{csr\_causal\_label\_audit.py}, \texttt{tr\_corpus\_audit.py})
  alt-süreç olarak çağırır; ek olarak render-crossref taraması, format
  ölçüm ve targets/renv durum kontrollerini kapsar (araç yoksa SKIP +
  neden).
\item
  Her madde için satır:
  \texttt{MADDE-ID\ \ DURUM(PASS/FAIL/SKIP/MANUEL)\ \ kısa\ kanıt}.
\item
  Sonda özet: toplam/PASS/FAIL/SKIP/MANUEL sayacı + genel sonuç.
\item
  \texttt{-\/-section\ \textless{}kod\textgreater{}} ile tek eksen;
  \texttt{-\/-chapter\ \textless{}dosya\textgreater{}} ile bölüm-düzeyi
  alt-küme; argümansız tam tez-düzeyi koşum.
\item
  Hiçbir dosyayı değiştirmez. Exit: 0 = otomatik maddelerin tümü
  PASS/SKIP; 1 = en az bir otomatik HARD FAIL.
\item
  Marmara format ölçümü (PDF varsa) ve render taraması ağır olabilir;
  bunlar \texttt{-\/-fast} ile atlanabilir (SKIP olarak işaretlenir).
\end{itemize}

\subsection{R5 --- Mevcut checklistleri
genişletme}\label{r5-mevcut-checklistleri-geniux15fletme}

\begin{itemize}
\tightlist
\item
  \texttt{format-kontrol-listesi.md},
  \texttt{kanit-ve-gizlilik-kontrol-listesi.md} ve
  \texttt{render-bagimliliklari.md} içinde master checklist'e ve
  script'e \textbf{çapraz bağlantı} eklenir (hangi madde script'in hangi
  kontrolüne karşılık geliyor).
\item
  Yeni eksenler (pipeline tazeliği, PDF format ölçümü) uygun mevcut
  dosyaya veya master belgeye eklenir; içerik tekrarı yerine referans
  tercih edilir.
\item
  \texttt{04\_kalite-kontrol/README.md}, yeni master checklist ve
  script'i akış içinde konumlandıracak biçimde güncellenir.
\end{itemize}

\subsection{R6 --- Araç-madde eşleme
tablosu}\label{r6-arauxe7-madde-eux15fleme-tablosu}

\begin{itemize}
\tightlist
\item
  Master belgede, her mevcut aracın (\texttt{scripts/util/*.py}, slash
  komutlar) hangi checklist maddelerini kapsadığını gösteren bir
  \textbf{kapsama matrisi} bulunur. Böylece hangi maddenin otomatik,
  hangisinin manuel kaldığı ve otomasyon boşlukları görünür olur.
\end{itemize}

\subsection{R7 --- Doğrulama (bu işin kendi
kontrolü)}\label{r7-doux11frulama-bu-iux15fin-kendi-kontroluxfc}

\begin{itemize}
\tightlist
\item
  \texttt{tez\_checklist\_verify.py} repo mevcut durumunda koşulur; hata
  vermeden (kod hatası anlamında) tamamlanır ve tutarlı bir rapor
  üretir.
\item
  Script'in raporladığı otomatik maddeler, elle koşulan tekil araç
  sonuçlarıyla tutarlı olmalı (ör. \texttt{bib\_hygiene} HARD sayısı
  script raporuyla aynı).
\item
  Master checklist'teki her OTOMATIK madde, script'te karşılık gelen bir
  kontrole sahip (yetim madde yok); her script kontrolü belgede bir
  maddeye bağlı.
\end{itemize}

\section{Kabul Kriterleri}\label{kabul-kriterleri}

\begin{itemize}
\tightlist
\item[$\square$]
  \texttt{tez-yazim/04\_kalite-kontrol/tez-kontrol-checklisti.md}
  mevcut; Türkçe; bölüm-düzeyi (A) + tez-düzeyi (B) iki katman içeriyor.
\item[$\square$]
  Belge R2'deki 8 eksenin tümünü madde madde kapsıyor.
\item[$\square$]
  Her madde R3 şemasındaki alanları (ID, mekanizma, komut, PASS ölçütü,
  tür, kanıt konumu, kapı/kaynak) taşıyor.
\item[$\square$]
  \texttt{scripts/util/tez\_checklist\_verify.py} mevcut; salt-okuma;
  argümansız tam koşum + \texttt{-\/-section} + \texttt{-\/-chapter} +
  \texttt{-\/-fast} destekli.
\item[$\square$]
  Script tüm otomatik araçları çağırıp birleşik PASS/FAIL/SKIP/MANUEL
  raporu ve özet exit kodu üretiyor; repo mevcut durumunda hatasız
  koşuyor.
\item[$\square$]
  Script çıktısı tekil araç sonuçlarıyla tutarlı (ör. bib\_hygiene HARD
  sayısı eşleşiyor).
\item[$\square$]
  Master belgede araç→madde kapsama matrisi var; her OTOMATIK madde
  script kontrolüyle, her script kontrolü bir maddeyle eşleşiyor (yetim
  yok).
\item[$\square$]
  Mevcut checklistler + \texttt{04\_kalite-kontrol/README.md} yeni
  belge/script'e çapraz-bağlı; içerik tekrarı yerine referans
  kullanılmış.
\item[$\square$]
  Değişiklikler tek toplu commit (Türkçe mesaj +
  \texttt{Co-authored-by:\ Ona\ \ \ \ \ \textless{}no-reply@ona.com\textgreater{}});
  yalnız QC katmanı dosyaları + yeni script.
\end{itemize}

\section{Uygulama Adımları
(sıralı)}\label{uygulama-adux131mlarux131-sux131ralux131}

\begin{enumerate}
\def\labelenumi{\arabic{enumi}.}
\tightlist
\item
  \textbf{Envanter sabitleme:} \texttt{scripts/util/} araçlarının her
  birinin gerçek arayüzünü (argümanlar, exit kodları, çıktı biçimi) ve
  slash komutların koştuğu komutları çıkar. Her aracın hangi eksen(ler)i
  kapsadığını netleştir (R6 matrisinin ham verisi).
\item
  \textbf{Eksen → madde ayrıştırması:} R2'deki 8 ekseni somut, atomik
  maddelere böl; her maddeye R3 şemasını doldur (ID, mekanizma, komut,
  PASS eşiği, tür, kanıt konumu, kapı/kaynak). Otomatik vs manuel
  ayrımını burada sabitle.
\item
  \textbf{Master checklist yaz (R1):} \texttt{tez-kontrol-checklisti.md}
  oluştur; (A) bölüm-düzeyi + (B) tez-düzeyi katmanlar; her eksen için
  madde tablosu + R6 kapsama matrisi + kanonik otorite referansları.
\item
  \textbf{Doğrulama script'i yaz (R4):}
  \texttt{tez\_checklist\_verify.py} --- mevcut araçları alt-süreç
  olarak çağıran orkestratör; render-crossref taraması, targets/renv
  durumu, (opsiyonel) PDF format ölçümü;
  \texttt{-\/-section/-\/-chapter/-\/-fast}; birleşik rapor + özet exit
  kodu; salt-okuma.
\item
  \textbf{Script--belge senkronu:} Her OTOMATIK maddenin script'te
  karşılığı, her script kontrolünün belgede maddesi olduğunu doğrula
  (yetim madde/kontrol yok).
\item
  \textbf{Mevcut checklistleri genişlet (R5):}
  \texttt{format-kontrol-listesi.md},
  \texttt{kanit-ve-gizlilik-kontrol-listesi.md},
  \texttt{render-bagimliliklari.md} ve \texttt{README.md}'ye
  çapraz-bağlantı + yeni eksen referansları ekle.
\item
  \textbf{Script'i koş ve doğrula (R7):}
  \texttt{tez\_checklist\_verify.py} repo mevcut durumunda koştur;
  raporun tutarlı ve hatasız üretildiğini, exit kodunun beklendiği gibi
  olduğunu, tekil araç sonuçlarıyla eşleştiğini teyit et.
\item
  \textbf{Teslim:} QC dosyaları + yeni script'i tek toplu commit ile ver
  (Türkçe mesaj + co-author).
\end{enumerate}

\section{Riskler / Notlar}\label{riskler-notlar}

\begin{itemize}
\tightlist
\item
  \textbf{Araç exit-kod tutarsızlığı:} Bazı araçlar (ör.
  \texttt{bib\_hygiene}) mevcut repo durumunda HARD fail (exit 1)
  döndürüyor olabilir. Script bunu doğru şekilde FAIL raporlar; bu,
  script'in bozuk olduğu anlamına gelmez --- checklist gerçek durumu
  yansıtır. Bu maddeler raporda net ayrılır.
\item
  \textbf{Ağır kontroller:} Tam \texttt{tar\_make()} ve PDF render
  dakikalar sürebilir; \texttt{-\/-fast} ile bunlar SKIP edilebilir,
  böylece hızlı ön-kontrol mümkün olur.
\item
  \textbf{Manuel maddeler kaçınılmaz:} Retorik akış, thick-description,
  jüri/kişisel alan doğruluğu gibi maddeler otomatikleştirilemez; belge
  bunları açıkça MANUEL işaretler ve script FAIL saymaz.
\item
  \textbf{PII sınırı:} Script salt-okuma ve satır-düzeyi veri sızdırmaz;
  \texttt{data/raw\textbar{}\ identified\textbar{}cleaned\textbar{}backup}
  içeriğini okumaz, yalnız varlık/gitignore durumunu denetler.
\item
  \textbf{Kapsam disiplini:} Checklist tespit eder; bulunan içerik
  eksiklerini düzeltmek ayrı iştir (bu iş yalnız kontrol altyapısını
  kurar).
\end{itemize}




\end{document}
